% ---------------------------------------------------------------------
%  5  CONCLUSION
% ---------------------------------------------------------------------
\section{Conclusion}
\label{sec:conclusion}

The 24 hourly price changes of Bitcoin's next day cannot be forecast
better than ``no change'': at the hourly scale the price is close to a
random walk, and only the next hour carries a small reversal signal (AUC
0.54). The size of the changes is predictable, so we recommend answering
the task with the expected size and an 80\% range for each hour. The
average of zero-shot Timestral and monthly-refitted Gradient Boosting
explains 16.4\% of the hourly variance (2.9 times a volatility rule), and
conformally calibrated ranges cover 80\% of the hours, release hours
included.

We learned that the evaluation design matters more than the choice of
model: chronological blocks with a purge, choices made without the test
block, comparison with simple rules and bootstrap intervals are what
separate a real signal from noise. Recency matters more than data volume
or extra sources: monthly refits were the largest single gain, while the
new features, ETH and macro variables added nothing reliable to the size
forecasts, and Binance order flow only a little to the next-hour
direction. A pretrained model, used as it comes, can match careful
feature engineering where there is signal, but cannot create signal where
there is none.

The study has limitations: a single test period; many tests on the return
target without a multiple-testing correction; no trading costs in the
next-hour signal; release hours that remain hard to forecast; and a
possible overlap between Timestral's training data and our test period.
Future work includes a rolling-origin evaluation over 2023--2026,
Holm or Benjamini--Hochberg corrections~\cite{benjamini1995}, a
cost-aware test of the next-hour signal, a forward test of Timestral on
data after its release, and a volatility-targeting design tuned on
validation.
